%Author1: Ganesh Ashok Satpute
%Webpage: ganeshasatpute.github.io
%Author2: Prof. Syed Abbas, IIT Mandi, India
%This file was used. https://math.mit.edu/~psh/exam/examdoc.pdf

%Author3: Jiayang Song
%This file is used to generate final exams for CS/SE110 Programming Language in MUST in Fall 2025

%%%%%%%%%%%%%%%%%%%%%%%%%%%%%%%%%%%%%%%%%%%%%%%%%%%%%%%%%%%%%%%%%%%%%%%%%%%%%%%%%%%%

% For question-only, use the following command.
% \documentclass[12pt, a4paper, addpoints]{exam}

% For answers, use the following command.
\documentclass[12pt, a4paper, answers, addpoints]{exam}


\def\CourseTitle{Computer Programming}
\def\ExamName{Final Examination}
\def\CourseNumber{CS110 \& SE110 \& BE104}
\def\ExamDate{2025/01/07}
\def\Duration{2}
\def\ExamTime{15:00 - 17:00}
\def\Term{2509}
\def\ExamDuration{2 Hours}


\usepackage[left=2cm, right=2cm, top=2cm, bottom=2cm]{geometry}

\usepackage{amsmath}
\usepackage{amsthm}
\usepackage{amsfonts}
\usepackage{amssymb}
\usepackage{mathrsfs}
\usepackage{tasks}
% \usepackage[margin=1in]{geometry}
\checkboxchar{$\Box$}
\usepackage{setspace}
%\usepackage{bbding}
%\usepackage{pifont}
\usepackage{wasysym}
\usepackage[shortlabels]{enumitem}
\renewcommand{\qedsymbol}{$\blacksquare$}
\usepackage{array} 
\usepackage{multicol}
\setlength{\columnsep}{1cm}
\usepackage{afterpage}
\usepackage{color} 
\usepackage{graphicx} % in preamble
\usepackage{listings}
\usepackage{upquote}

\lstset{
    escapeinside={(*@}{@*)},           % Allow LaTeX commands between (*@ and @*)
}

% new macros
\newcommand{\todo}[1]{\textcolor{red}{(To do): #1}}
\newcommand{\tored}[1]{\textcolor{red}{\textbf{#1}}}
\usepackage{xeCJK}
\defaultfontfeatures{Ligatures=TeX}

\setCJKmainfont[
    BoldFont={Noto Sans TC}
]{Noto Sans TC}

% Required by listings/typewriter text
\setCJKmonofont{Noto Sans TC}

\newcommand\blankpage{%
    \null
    \thispagestyle{empty}%
    \addtocounter{page}{-1}%
    \shownoanswers{\newpage}}
%% https://tex.stackexchange.com/a/203986
% Create a True-False question format
\newcommand*{\TrueFalse}[1]{%
\ifprintanswers
    \ifthenelse{\equal{#1}{T}}{%
        \textbf{TRUE}\hspace*{14pt}False
    }{
        True\hspace*{14pt}\textbf{FALSE}
    }
\else
    {True}\hspace*{20pt}False
\fi
} 
%% The following code is based on an answer by Gonzalo Medina
%% https://tex.stackexchange.com/a/13106/39194
\newlength\TFlengthA
\newlength\TFlengthB
\settowidth\TFlengthA{\hspace*{1.16in}}
\newcommand\TFQuestion[2]{%
    \setlength\TFlengthB{\linewidth}
    \addtolength\TFlengthB{-\TFlengthA}
    \parbox[t]{\TFlengthA}{\TrueFalse{#1}}\parbox[t]{\TFlengthB}{#2}}
%%

%%
\pagestyle{headandfoot}
%\extraheadheight{.25in}
\runningheader{{\CourseNumber} - {\CourseTitle}}
{}
{{\ExamName} - {\Term}} 
\firstpagefooter{}{Page \thepage\ of \numpages}{Please go on to the next page\ldots} %
\runningfooter{}{Page \thepage\ of \numpages}{\iflastpage{End of exam.}{Please go on to the next page\ldots}}

\renewcommand{\fillinlinelength}{1cm}  % change default length

\renewenvironment{solution}
  {\par\noindent\textbf{Solution: }\color{red}} % begin code
  {\par}    

% --- show solution chunks only when 'answers' is active ---
\usepackage{xcolor} % for \color

% Inline macro: use as \showsol{<solution text>}
\newcommand{\showsol}[1]{\ifprintanswers{\textbf{\color{red}\underline{Solution:} }\color{red}#1}\fi}

\newenvironment{show_code_sol}
  {\par\noindent\textbf{Solution:}\par
   \lstset{basicstyle=\ttfamily\small\color{red}}}
  {\par}


\usepackage{listings}
\usepackage{comment} % for \excludecomment (safe with verbatim)
% \usepackage[T1]{fontenc}


% Environment version that mimics 'solution' but guarded
\newenvironment{solonly}
  {\ifprintanswers\par\noindent\textbf{Solution: }\color{red}}
  {\par\fi}

% If you want to keep using the built-in 'solution' name,
% make sure it stays conditional:
\renewenvironment{solution}
  {\ifprintanswers\par\noindent\textbf{Solution: }\color{red}}
  {\par\fi}

\newcommand{\shownoanswers}[1]{\ifprintanswers\else #1\fi}

% Base style for C
% One listings + comment load
\usepackage{listings}
\usepackage{comment} % safe with verbatim

% --- styles (no frame, no line numbers) ---
\lstdefinestyle{examc}{
  language=C,
  basicstyle=\ttfamily\small,
  numbers=none,
  frame=none,
  columns=fullflexible,
  keepspaces=true,
  showstringspaces=false,
  tabsize=4,
  breaklines=true,
  keywordstyle=\bfseries,
  commentstyle=\itshape,
  stringstyle=\ttfamily
}
\lstdefinestyle{examc-sol}{
  style=examc,
  basicstyle=\ttfamily\small\color{red}
}

% --- Environments ---
\lstnewenvironment{Ccode}[1][]
  {\lstset{style=examc,#1}}
  {}

% Define CcodeSol only when answers are printed; otherwise exclude it.
\ifprintanswers
  \lstnewenvironment{CcodeSol}[1][]
    {\lstset{style=examc-sol,#1}}
    {}
\else
  \excludecomment{CcodeSol} % body is ignored; no verbatim processing
\fi

% Inline C
\newcommand{\Cinline}[1]{\lstinline[language=C]{#1}}

% ---------------------------------------------------------------------
  


%%%%%%%%%%%%%%%%%%%%%%%%%%%%%%%%%%%%%%%%%%%%%%%%%%%%%%%%%%%%%%%%%%%%%%%%%%%%%%%%%%%%
\begin{document}


\noindent
\begin{minipage}[t]{0.48\textwidth}
    \vspace{-10pt}
    \fbox{\parbox[c][1.2cm][c]{4cm}{Seat\\No.}}
\end{minipage}
\hfill
% \begin{minipage}[t]{0.48\textwidth}
%     \vspace{-10pt}
%     \centering
%     \includegraphics[width=0.35\textwidth]{must-logo.pdf}
% \end{minipage}

\vspace{-2cm} % adjust as needed
\begin{center}
    \includegraphics[width=0.27\textwidth]{must-logo.pdf}
\end{center}



\begin{center}
    
    \vspace*{0.5cm}    
    {\textbf{\Huge 澳\hspace{2em}門\hspace{2em}科\hspace{2em}技\hspace{2em}大\hspace{2em}學}}\\[0.3cm]
    {\large \textbf{MACAU UNIVERSITY OF SCIENCE AND TECHNOLOGY}}\\[1cm]

    {\Large \textbf{School of Computer Science and Engineering}}\\[0.2cm]
    {\Large \textbf{Faculty of Innovation Engineering}}\\[0.2cm]
    {\Large \textbf{Bachelor of Science}}\\[1cm]

    {\large \textbf{2025/2026 – 1st Semester Final Question Booklet}}\\[1cm]

    {\Large \textbf{\CourseNumber \quad \CourseTitle}}\\[1cm]
\end{center}

% \noindent
% \hspace{2cm}% adjust indent width here
\begin{center}
\begin{tabular}{@{}p{4cm}p{8cm}@{}} 
    \textbf{No. of Pages:} & {\numpages} Pages (including this page)\\[0.5cm]
    \textbf{Exam Date:} & \ExamDate  \\[0.5cm]
    \textbf{Exam Duration:} & {\ExamDuration} ({\ExamTime})\\[0.5cm]
    \textbf{Student Name:} & \underline{\hspace{7cm}}\\[0.5cm]
    \textbf{Student Number:} & \underline{\hspace{7cm}}\\
\end{tabular}
\end{center}

\vspace{0.5cm}

\noindent
\textbf{Please read the following instructions carefully.}
\begin{itemize}[leftmargin=*]
    \item This question booklet contains \textbf{\numquestions} questions, \textbf{\numpages} pages (including the cover) for the total of \textbf{\numpoints} points/marks. Check whether all pages are present.
    % Please write your answers on the Question Booklet.
    
    \item Return this Question Booklet, together with all extra unused answer sheets and draft paper, at the end of the examination.
    
    \item This is a \textbf{Closed Book Exam}. The use of dictionaries, calculators, or electronic devices is \textbf{NOT permitted}.
    
    \item \textbf{Answers must be provided in English to receive marks; failure to do so will result in no marks being awarded.}

    % \item \textbf{Read the instructions for individual questions carefully} before answering the questions. Clearly identify your final answers.
\end{itemize}


\vfill

\hspace{0.5in}
\begin{questions}
\addpoints


\newpage


\question[2]{When we write functions in C, we often need to pass data or variables from one function to another. Define what is meant by \emph{call by value} in C.
Why can't a function directly modify a variable passed from \texttt {main()}?}

\shownoanswers{\vspace{6cm}}

\showsol{
C uses \emph{call by value}: the function receives a \emph{copy} of each argument. 
Therefore, it cannot directly modify the original variable in \texttt{main()} because changes affect only the local copy. 

\textbf{Remark: Give full marks if the student mentions that arguments are passed as \emph{copies}.}
\vspace{1cm}
}



\question[4]{Consider the code below:
\vspace{-20pt}
\begin{center}
\begin{lstlisting}[basicstyle=\small\ttfamily]
        #include <stdio.h>
        int main(void) {
            int arr[2][3] = {{1 ,2 ,3}, {4, 5, 6}};
            int (*p)[3] = arr;
            return 0;
        }
\end{lstlisting}
\end{center}
% What does \texttt{*(arr + 3)} evaluate to?} Explain how pointer arithmetic works in this example.

\begin{parts}
    \part (2 points) In most expressions, what is the data type of \texttt{arr}? What is the data type of \texttt{*arr}? 
    \shownoanswers{\vspace{5cm}}\\
    \showsol{
        The data type of \texttt{arr} is: \textbf{\texttt{int (*)[3]}} or a pointer to an array of three integers.
        The data type of \textbf{\texttt{*arr}} is: \texttt{int *} or a pointer to an integer.\\
        \textbf{Remark: 1 point for each.}
    \vspace{0.5cm}
    }
    
    \part (2 points) How to represent the element \texttt{arr[1][2]} with pointer arithmetic (using \texttt{p})?
    \shownoanswers{\vspace{3cm}}\\
        \showsol{
        \texttt{*((*p+1)+2)}
    \vspace{0.5cm}
    }

\end{parts}
}

% \shownoanswers{\vspace{4cm}}

% \showsol{
% \texttt{*(arr + 3)} evaluates to \textbf{40}. \\
% \texttt{arr} points to \texttt{arr[0]}. Adding 3 advances by \texttt{3 * sizeof(int)} bytes to \texttt{arr[3]}; dereferencing yields \texttt{arr[3]} which is 40.

% \textbf{Remark: One point for \textbf{40}, one point for the explanation.}
% \vspace{1cm}
% }




% \question[2] {Explain the difference between the character constant \texttt{'A'} and the string constant \texttt{"A"}.}

% \shownoanswers{\vspace{4cm}}

% \showsol{
% \texttt{'A'} is a \emph{character constant} of type \texttt{int} with value 65 (ASCII for 'A'). \\
% \texttt{"A"} is a \emph{string literal}: an array of two chars \{\texttt{'A'}, \texttt{'\textbackslash0'}\}. 

% \textbf{Remark: Give full marks if the student notes that \texttt{"A"} includes an additional null terminator \texttt{'\textbackslash0'}.}

% \vspace{1cm}
% }

\shownoanswers{\newpage}
\question[2]{Are there any possible bugs in the code below? If yes, explain your answer and provide the corrected version of the code.
\vspace{-20pt}
\begin{center}
\begin{lstlisting}[basicstyle=\small\ttfamily]
        #include <stdio.h>
        int main(void) {
            int arr[3] = {10, 20, 30};
            for (int i = 0; i <= 3; i++)
                printf("%d\n", arr[i]);
            return 0;
        }
\end{lstlisting}
\end{center}
}

\shownoanswers{\vspace{6cm}}

\showsol{
Potential bug: array out-of-bounds access. 
The array arr has only three elements with valid indices 0, 1, and 2.
However, the loop condition i <= 3 allows i to reach the value 3, leading to an attempt to access arr[3], which is out of bounds. \\
\textbf{Fix:} the loop condition should be changed to \texttt{for (int i = 0; i < 3; i++)} or \texttt{for (int i = 0; i <= 2; i++)}

\textbf{Remark: Give full marks if the student mentions array out-of-bounds access}
\vspace{1cm}
}


% \shownoanswers{\newpage}
\question[2]{Explain the difference between \texttt{break} and \texttt{continue} in loop control statements.}

\shownoanswers{\vspace{7cm}}

\showsol{
\texttt{break} exits the \emph{nearest} loop immediately. \\
\texttt{continue} skips the rest of the current loop iteration and proceeds to the next iteration (re-checking the condition).

\textbf{Remark: Students must mention the effect of each statement on the loop or iteration. Reflect CO5}
}



\question[2]{Explain the difference between \textbf{internal linkage} and \textbf{external linkage} in C.}

\shownoanswers{\vspace{5cm}}

\showsol{
\textbf{External linkage} means the identifier can be accessed across multiple files (translation units).\\
\textbf{Internal linkage} means the identifier is restricted to the current file only.  


\textbf{Remark: One point for each explanation. It is fine if students mention a variable instead of an identifier.}

\vspace{1cm}
}

\shownoanswers{\newpage}
\question[2]{Explain the difference between a \texttt{struct} and an \texttt{array}.}

\shownoanswers{\vspace{6cm}}

\showsol{
A \texttt{struct} groups \emph{heterogeneous}, named fields that can have different types; elements are accessed by field name. 
An \texttt{array} is a contiguous sequence of \emph{homogeneous} elements of the same type, accessed by index.

\textbf{Remark: Give full marks if the student mentions that a \texttt{struct} can hold data of different types, while an \texttt{array} cannot. Reflect CO2.}
\vspace{1cm}
}

\question[2]{What happens if you open a file with \texttt{"w"} mode using \texttt{fopen()} and the file already exists?}

\shownoanswers{\vspace{6cm}}

\showsol{
Opening with mode \texttt{"w"} \emph{truncates} the file to zero length (data lost) or creates a new file if it doesn't exist. 

\textbf{Remark: full marks if the student mentions that all original file contents are overwritten (truncated).}
\vspace{1cm}
}

\question[2]{Consider the following C program that uses command-line arguments:

\begin{lstlisting}[basicstyle=\small\ttfamily]
        #include <stdio.h>
        int main(int argc, char *argv[]) {
            printf("argc = %d\n", argc);
            printf("argv [1] = %s\n", argv[1]);
            return 0;
        }
\end{lstlisting}
}

% \shownoanswers{\vspace{4cm}}

\begin{parts}
    \part (1 point) What does \texttt{argc} represent?  
    \shownoanswers{\vspace{3cm}}\\
    \showsol{
        \texttt{argc} is the number of command-line arguments, including the program name.  \\
        \textbf{Remark: Reflect CO4.}
    \vspace{0.5cm}
    }
    
    \part (1 point) What does \texttt{argv[1]} represent?  
    % \shownoanswers{\vspace{3cm}}\\
        \showsol{
        \texttt{argv[1]} is the first argument provided by the user (a string).  \\
        \textbf{Remark: Reflect CO2.}
    \vspace{0.5cm}
    }

\end{parts}

\shownoanswers{\newpage}
\question[3]{What happens in this example? (Explain why the result might be negative.)
\vspace{-20pt}
\begin{center}
    \begin{lstlisting}[basicstyle=\small\ttfamily]
    #include <stdio.h>
    int main(void) {
            unsigned int a = 3E9;
            int b = (int)a;
            printf("%d", b);
        return 0;
    }
    \end{lstlisting}
\end{center}
}


\shownoanswers{\vspace{6cm}}

\showsol{
\texttt{3E9} is a \texttt{double} (3{,}000{,}000{,}000).
Casting it to \texttt{int} may cause overflow (implementation-defined).  
On a typical 32-bit two’s complement system, it wraps around to \textbf{-1294967296}, which explains the negative output.

\textbf{Remark: Give full marks if the students mention type casting causing overflow. Reflect CO1.}
\vspace{1cm}
}

\question[3] {What will the following code print? (Show the full output) 
\vspace{-20pt}
\begin{center}
    \begin{lstlisting}[basicstyle=\small\ttfamily]
    #include <stdio.h>
    int main(void) {
        int i = 0;
        while (i++ < 5)
            printf("%d ", ++i);
        return 0;
    }
    \end{lstlisting}
\end{center}
}

\shownoanswers{\vspace{3cm}}

\showsol{
\textbf{2, 4, 6}.  \\
 Each loop test uses \texttt{i++} (post-increment), then the body prints \texttt{++i} (pre-increment), producing 2, 4, 6 before the condition fails.

\textbf{Remark: one point for each correctly printed number. Reflect CO3.}
\vspace{1cm}
}


\shownoanswers{\newpage}
\question[4]{In C, what is the purpose of \texttt{malloc()} and \texttt{free()}?  Identify one potential problem in the following code:
\vspace{-20pt}
\begin{center}
    \begin{lstlisting}[basicstyle=\small\ttfamily]
        #include <stdio.h>
        #include <stdlib.h>
        
        int main() {
            int *p = malloc(5 * sizeof(int));
            for (int i = 0; i < 5; i++)
                p[i] = i * 2;
            free(p);
            printf("%d\n", p[0]);
            return 0;
        }
    \end{lstlisting}
\end{center}
}

\shownoanswers{\vspace{6cm}}

\showsol{
\texttt{malloc()} dynamically allocates memory from the heap; \texttt{free()} releases that memory.  
% In this code, 5 integers are allocated and initialized, then the memory is freed.

\textbf{The problem}: after calling \texttt{free(p)}, the pointer becomes a \emph{dangling pointer}.  
Accessing \texttt{p[0]} after freeing the memory is \textbf{undefined behavior} (could crash or print garbage).

\textbf{Remark: 2 points for explaining the purposes of \texttt{malloc()} and \texttt{free()}. 2 points for identifying and explaining the problem.}

\vspace{1.5cm}
}

\question[4]{Consider the following C program:
\vspace{-20pt}
\begin{center}
    \begin{lstlisting}[basicstyle=\small\ttfamily]
    #include <stdio.h>
    #include <string.h>
    
    int main(void) {
        char str[20] = "C Programming";
        char *p = str;
        while (*p) {
            if (*p == 'a')
                *p = '@';
            p++;
        }
        printf("%s\n", str);
        return 0;
    }
    \end{lstlisting}
\end{center}

\begin{parts}
    \part (2 points) What will be printed on the screen?
    \shownoanswers{\vspace{3cm}}\\
    \showsol{
        \texttt{C Progr@mming}\\
        \textbf{Remark: Reflect CO4.}
    \vspace{0.5cm}
    }
    
    \part (2 point) How many characters are in the string \texttt{str} (excluding \texttt{'\textbackslash0'})
    \shownoanswers{\vspace{3cm}}\\
        \showsol{
        13 characters (excluding \texttt{'\textbackslash0'}). \\
        \textbf{Remark: Reflect CO2.}
    \vspace{0.5cm}
    }

\end{parts}
}

% \shownoanswers{\newpage}
\question[4]{Write the exact output of the following program, including spaces and newlines if any.
\vspace{-20pt}
\begin{center}
    \begin{lstlisting}[basicstyle=\small\ttfamily]
        #include <stdio.h>
        void rev(const char *s) {
            if (*s == '\0') 
                return;
            rev(s + 1);
            putchar(*s);
        }
            
        int main(void) {
            const char *t = "ab cd";
            rev(t);
            putchar('\n');
            return 0;
        }
    \end{lstlisting}
\end{center}
}

\shownoanswers{\vspace{4cm}}

\showsol{
\texttt{dc ba}\\
\textbf{Remark: The string is printed in reverse order, preserving the space.} 
\vspace{1cm}
}

\question[4]{Write the exact output of the following program, including spaces if any.
\vspace{-20pt}
\begin{center}
    \begin{lstlisting}[basicstyle=\small\ttfamily]
#include <stdio.h>
int main(void) {
    int arr[4] = {7, 8, 9, 10};
    int *ap[2] = { &arr[1], &arr[3] }; 
    int (*pa)[4] = &arr; 
    
    printf("%d %d %d %d\n", *ap[0], *ap[1], (*pa)[0], *(*pa + 2));
    return 0;
}
    \end{lstlisting}
\end{center}
}

\shownoanswers{\vspace{3cm}}

\showsol{
\texttt{8 10 7 9}\\
\texttt{*ap[0] = arr[1] = 8}, \texttt{*ap[1] = arr[3] = 10}, \texttt{(*pa)[0] = arr[0] = 7}, \texttt{*(*pa + 2) = arr[2] = 9}.

\textbf{Remark: one point for each correctly identified output value. Reflect CO3.}

\vspace{1.5cm}
}

\shownoanswers{\newpage}
\question[5]{Fill the blanks so the program reads a string and prints its length using the function \texttt{mystrlen()}.

\textbf{Example Input:} \texttt{Hello} \\
\textbf{Example Output:} \texttt{Length = 5}
\vspace{-10pt}
\begin{center}
    \begin{lstlisting}[basicstyle=\small\ttfamily]
    #include <stdio.h>
    
    int mystrlen(char s[]) {
        int i = 0;
        while (s[(*@\fillin[\tored{i}]@*)] != (*@\fillin[\tored{'\textbackslash0'}]@*)) // fill blanks
            i++;
        return (*@\fillin[\tored{i}]@*);            // fill blank
    }
    
    int main(void) {
        char str[100];
        scanf("%s", (*@\fillin[\tored{str}]@*));    // fill blank
        printf("Length = %d\n", mystrlen((*@\fillin[\tored{str}]@*))); // fill blank
        return 0;
    }
    \end{lstlisting}
\end{center}
}
\shownoanswers{\vspace{2cm}}

\showsol{    
    \textbf{Remark: 1 point for each correctly filled blank; explanations are not required. Reflect CO2.}
    \vspace{1cm}
}


\question[6] {Complete the blanks so the program reads numbers from a file (\texttt{numbers.txt}) and prints the total sum.

Assume the file \texttt{numbers.txt} contains content: \texttt{1 2 3 4 5}\\
\textbf{Expected Output:} \texttt{Sum = 15}
\vspace{-10pt}
\begin{center}
    
    \begin{lstlisting}[basicstyle=\small\ttfamily]
    #include <stdio.h>
    
    int main(void) {
        FILE *fp;
        int x, sum = 0;
        fp = fopen((*@\fillin[\tored{"numbers.txt"}][2.5cm]@*), (*@\fillin[\tored{"r"}]@*));           // fill blanks
        if (fp == NULL) {
            printf("Cannot open file.\n");
            return 1;
        }
        
        while (fscanf((*@\fillin[\tored{fp}]@*), (*@\fillin[\tored{"\%d"}]@*), (*@\fillin[\tored{\&x}]@*)) == 1) {  // fill blanks
            sum += x;                 
        }
        
        fclose((*@\fillin[\tored{fp}]@*));                       // fill blank
        printf("Sum = %d\n", sum); 
        return 0;
    }
    \end{lstlisting}
\end{center}

}
\shownoanswers{\vspace{2cm}}

\showsol{    
    \textbf{Remark: 1 point for each correctly filled blank; explanations are not required. Reflect CO4.}
    \vspace{1cm}
}



\question[4] {Fill the blanks so the program defines a structure using \texttt{typedef} and accesses its members via a pointer.

\textbf{Expected Output:} \texttt{Car: Tesla, Year: 2024}
\vspace{-10pt}

\begin{center}
\begin{lstlisting}[basicstyle=\small\ttfamily]
#include <stdio.h>

typedef struct {
    char brand[20];
    int year;
} (*@\fillin[\tored{Car}]@*);       // fill blank

int main(void) {
    Car c1 = {"Tesla", 2024};       
    Car *p = (*@\fillin[\tored{\&c1}]@*);        // fill blank
    printf("Car: %s, Year: %d\n", (*@\fillin[\tored{p->brand}][1.5cm]@*), (*@\fillin[\tored{p->year}][1.5cm]@*)); // fill blanks
    return 0;
}
\end{lstlisting}
\end{center}
}
\shownoanswers{\vspace{2cm}}

\showsol{    
    \textbf{Remark: 1 point for each correctly filled blank; explanations are not required.}
    \vspace{1cm}
}


\shownoanswers{\newpage}
\question[10] {Write a C program that computes the transpose of a matrix (2D array). Partial code is given; you need to complete this program. The transpose of a matrix is a new matrix obtained by swapping its rows and columns. For example:

The transpose of a matrix
 $\begin{bmatrix}
  1 & 2 & 3\\ 
  4 & 5 & 6\\
\end{bmatrix}$ is 
 $\begin{bmatrix}
  1 & 4\\ 
  2 & 5\\
  3 & 6
\end{bmatrix}$.
The transpose of a matrix
 $\begin{bmatrix}
  1 & 2 \\ 
  3 & 4 \\
\end{bmatrix}$ is 
 $\begin{bmatrix}
  1 & 3\\ 
  2 & 4\\
\end{bmatrix}$.


\textbf{Expected Output:}  
$\begin{matrix}
    4 & 9 & 7 \\
    2 & 3 & 4 
\end{matrix}$ 

\textbf{Program to Complete:}
\begin{lstlisting}[basicstyle=\small\ttfamily]
#include <stdio.h>
int main(void) {

    int arr[3][2] = {{4, 2}, {9, 3}, {7, 4}};
    // Complete the program 
    
 }
\end{lstlisting}

}

\begin{CcodeSol}
#include <stdio.h>
int main(void) {

    int arr[3][2] = {{4, 2}, {9, 3}, {7, 4}};
    int r = 3, c = 2;
    
     for (int j = 0; j < c; j++) {
         for (int i = 0; i < r; i++) {
             printf("%d ", arr[i][j]);
         }
         putchar('\n');
     }
     return 0;
 }
\end{CcodeSol}

\showsol{    
    The following marking guidance is for reference only; you can make your own judgment. 
    \begin{itemize}
        \item If students have syntax errors, deduct up to 4 points.
        \item Give 2 points if the student demonstrates the correct idea of using nested \texttt{for} loops, even if minor errors are present.
        \item Give 1 point if the student correctly prints array elements using \texttt{printf("\%d ", arr[i][j]);}.
        \item Give 1 point if the student uses \texttt{putchar('\textbackslash n')} (or an equivalent method) to print each row on a new line.
    \end{itemize}
    \vspace{1cm}
}

\shownoanswers{\newpage}
\question[10] {Write a C program that sorts an array in ascending order and prints the sorted sequence on one line. Partial code is given; you only need to complete the \texttt{sort} function.

\begin{itemize}[leftmargin=*]
    
    % \item You must complete the \texttt{sort} function using pointer arithmetic only (no array indexing with \texttt{arr[]} anywhere in \texttt{sort} function).
    
    \item The program calls the \texttt{sort} function to sort the array and print the elements in ascending order.
    
\end{itemize}

\textbf{Expected Output:} \texttt{The sorted array is: -8 -5 1 3 4 14 }

\textbf{Program to Complete:}

\begin{lstlisting}[basicstyle=\small\ttfamily]
#include <stdio.h>
void sort(int *, int);

int main(void) {
    int arr[] = {4, 3, -5, 14, -8, 1};
    int N = sizeof(arr)/sizeof(arr[0]);

    sort(arr, N); 
    
    printf("The sorted array is: ");
    for (int i = 0; i < N; i++) {
        printf("%d ", arr[i]);
    }
    putchar('\n');
    return 0;
}

void sort(int *p, int n) {
    // Complete this function
    
}

\end{lstlisting}

}

\begin{CcodeSol}
#include <stdio.h>
void sort(int *, int);

int main(void) {
    int arr[] = {4, 3, -5, 14, -8, 1} ;
    int N = sizeof(arr)/sizeof(arr[0]);

    sort(arr, N);  // sort array in ascending order
    
    printf("The sorted array is: ");
    for (int i = 0; i < N; i++) {
        printf("%d ", arr[i]);
    }
    putchar('\n');
    return 0;
}

void sort(int *p, int n) {
    for (int i = 0; i < n - 1; i++) {
        for (int j = 0; j < n - 1 - i; j++) {
            if (p[j] > p[j + 1]) {
                int temp = p[j];
                p[j] = p[j + 1];
                p[j + 1] = temp;
            }
        }
    }
}

\end{CcodeSol}

\showsol{    
    The following marking guidance is for reference only; you can make your own judgment. 
    The solution provided is a sample bubble sort approach. There are various sorting algorithms; the answers do not need to be the same as the reference.
    \begin{itemize}
        \item Deduct up to 4 points for syntax errors.
        \item Give 4 points if the student demonstrates the idea of using nested \texttt{for} loops.
    \end{itemize}
    \vspace{1cm}
}


\shownoanswers{\newpage}
\question[10] {Write a C program that defines a structure called \texttt{Student} to store each student's name and score.  
You must write a function that computes and returns the \textbf{highest score} in the array of structures, and prints the \textbf{name and score} of the top student.
The initial code is given; complete this program.

\textbf{Expected Output:}

\texttt{Top Student: Carol (95)}


\textbf{Program to Complete:}

\begin{lstlisting}[basicstyle=\small\ttfamily]
#include <stdio.h>
#include <string.h>

struct Student {
    char name[50];
    int score;
};

int find_top(const struct Student *, int);

int main(void) {
    struct Student s[] = {
        {"Carol", 95},
        {"Alice", 90},
        {"Bob", 85}
    };
    int N = sizeof(s) / sizeof(s[0]);

    int index = find_top(s, N);
    printf("Top Student: %s (%d)\n", s[index].name, s[index].score);

    return 0;
}

int find_top(const struct Student *p, int n) {
    // Complete this function
}
\end{lstlisting}

}


\begin{CcodeSol}

#include <stdio.h>
#include <string.h>

struct Student {
    char name[50];
    int score;
};

int find_top(const struct Student *, int);

int main(void) {
    struct Student s[] = {
        {"Carol", 95},
        {"Alice", 90},
        {"Bob", 85}
    };
    int N = sizeof(s) / sizeof(s[0]);

    int index = find_top(s, N);
    printf("Top Student: %s (%d)\n", s[index].name, s[index].score);

    return 0;
}
int find_top(const struct Student *p, int n) {
    int max_index = 0;              // assume first student is top
    for (int i = 1; i < n; i++) {
        if (p[i].score > p[max_index].score) {
            max_index = i;
        }
    }
    return max_index;
}

\end{CcodeSol}

\showsol{    
    The following marking guidance is for reference only; you can make your own judgment. Reflect CO2.
    \begin{itemize}
        \item If students have syntax errors, deduct up to 4 points.
        \item Give 2 points if the students show the idea of using nested \texttt{for loops}, even if the loop may have mistakes.
        \item Give 2 points if the students know the idea of keeping track of the index with the highest score, and return that index to \texttt{main()}.
    \end{itemize}
    \vspace{1cm}
}

\end{questions}


\vfill
\noindent\makebox[2.5in]{\textbf{\Large Score:} \hrulefill}

\end{document}


